\documentclass[12pt,a4paper]{article}

\usepackage[spanish,es-nodecimaldot]{babel}
\usepackage[utf8]{inputenc}
\usepackage[T1]{fontenc}
\usepackage{lmodern}
\usepackage{geometry}
\usepackage{microtype}
\usepackage{setspace}
\usepackage{parskip}
\usepackage{enumitem}
\usepackage{xcolor}
\usepackage{listings}
\usepackage{hyperref}
\usepackage{amsmath}
\usepackage{amssymb}
\usepackage{array}
\usepackage{booktabs}
\usepackage{longtable}
\usepackage{tcolorbox}

\geometry{margin=2.5cm}
\onehalfspacing
\setlist[itemize]{leftmargin=*,itemsep=0.2em}
\setlist[enumerate]{leftmargin=*,itemsep=0.2em}

\hypersetup{
  colorlinks=true,
  linkcolor=black,
  urlcolor=blue,
  pdftitle={Resumen teorico TP8 - Mocks, EvoSuite, Randoop y ejecucion simbolica},
  pdfauthor={}
}

\newtcolorbox{definicionbox}[1][]{
  colback=blue!4!white,
  colframe=blue!50!black,
  fonttitle=\bfseries,
  title=Definicion,
  #1
}

\newtcolorbox{ideabox}[1][]{
  colback=green!4!white,
  colframe=green!40!black,
  fonttitle=\bfseries,
  title=Idea clave,
  #1
}

\newtcolorbox{ejemplobox}[1][]{
  colback=orange!4!white,
  colframe=orange!50!black,
  fonttitle=\bfseries,
  title=Ejemplo,
  #1
}

\lstdefinestyle{codestyle}{
  basicstyle=\ttfamily\small,
  frame=single,
  breaklines=true,
  showstringspaces=false
}

\title{\textbf{Resumen Teorico -- Practico 8}\\
Dobles de prueba, generacion automatica y ejecucion simbolica\\
\textit{Testing de Software 2024}}
\author{}
\date{}

\begin{document}

\maketitle
\tableofcontents
\newpage

\section*{Introduccion}
\addcontentsline{toc}{section}{Introduccion}

Este resumen sintetiza el material teorico asociado al TP8:

\begin{itemize}
  \item \textit{Notas 14: Dobles de pruebas (Mocking)} (V. Bengolea).
  \item \textit{Notas 15: Generacion automatica de tests --- EvoSuite y algoritmos
  evolutivos} (V. Bengolea).
  \item \textit{Notas 16: Ejecucion simbolica dinamica} (V. Bengolea).
\end{itemize}

\textbf{Tres ideas centrales que se cruzan.}

\begin{itemize}
  \item \textbf{Dobles de prueba}: como aislar la clase bajo test de sus colaboradores cuando
  estos son lentos, indeterministas o todavia no estan implementados.
  \item \textbf{Generacion automatica de tests}: dos enfoques representativos --- \emph{Randoop}
  (generacion aleatoria guiada por feedback) y \emph{EvoSuite} (busqueda evolutiva guiada por
  cobertura).
  \item \textbf{Ejecucion simbolica}: una tecnica que recorre el programa con \emph{valores
  simbolicos} en lugar de concretos, recolecta \emph{condiciones de camino} y las resuelve
  con un \emph{constraint solver} para fabricar inputs que recorran caminos especificos.
\end{itemize}

Las tres ideas se complementan: los mocks resuelven dependencias estaticas, la generacion
automatica produce volumen de tests, y la ejecucion simbolica garantiza precision en la
exploracion de caminos dificiles.

\section{Dobles de prueba (Mocking)}

\subsection{Por que se necesitan}

Cuando una unidad bajo prueba (SUT, \emph{System Under Test}) depende de otros componentes
(\emph{Depended-On Components}, DOCs), aparecen dos problemas:

\begin{itemize}
  \item Los DOCs pueden ser \emph{lentos} (red, base de datos), \emph{indeterministas}
  (reloj, IO), o \emph{costosos} (servicios pagos).
  \item Los DOCs pueden no estar \emph{implementados todavia}, especialmente en desarrollo
  agil.
\end{itemize}

\begin{definicionbox}
Un \textbf{doble de prueba} (\emph{test double}) es un esqueleto o implementacion especial de
un modulo que se necesita para testear un componente que depende de el. Reemplaza al DOC para
permitir que el SUT pueda probarse de manera aislada.
\end{definicionbox}

\subsection{Entrada y salida indirectas}

El SUT interactua con sus DOCs de dos maneras:

\begin{itemize}
  \item \textbf{Entrada indirecta}: el SUT \emph{lee} un valor del DOC (e.g.\
  \texttt{doc.meth(...)} devuelve algo que el SUT usa en su logica).
  \item \textbf{Salida indirecta}: el SUT \emph{modifica} el estado del DOC (e.g.\
  \texttt{db.save(o)}).
\end{itemize}

\begin{ideabox}
Un doble reemplaza al DOC para \emph{controlar las entradas indirectas} (forzar que el SUT
reciba el valor que necesitamos) y/o para \emph{monitorear las salidas indirectas}
(verificar que el SUT invoco al DOC correctamente).
\end{ideabox}

\subsection{Estrategias de integracion}

\paragraph{Adentro hacia afuera (inside-out).} Se codifica y testea primero el nucleo y se va
hacia las capas externas. Minimiza la necesidad de dobles pero minimiza tambien el feedback
temprano. Si los requisitos cambian, hay que retrabajar tests.

\paragraph{Afuera hacia adentro (outside-in).} Se codifica y testea desde las capas externas
hacia el nucleo, usando dobles que se van reemplazando a medida que se completa la
implementacion. Maximiza dobles y feedback, pero minimiza el trabajo adicional ante cambios
de requisitos. Es el patron natural en desarrollo agil.

\subsection{Taxonomia de dobles}

Hay cinco tipos principales (segun Gerard Meszaros, \emph{xUnit Test Patterns}):

\begin{enumerate}
  \item \textbf{Dummy}: solo satisface dependencias formales. Se pasa como parametro pero no
  hace nada; en general lanza excepcion si se lo usa. Util cuando un constructor requiere un
  parametro que el escenario no ejercita.
  \item \textbf{Stub}: provee respuestas \emph{prefabricadas} para las invocaciones que el
  SUT le hace. Controla las entradas indirectas.
  \item \textbf{Spy}: como un stub, pero ademas \emph{registra} las invocaciones recibidas
  para inspeccionarlas despues. Tambien permite monitorear salidas indirectas.
  \item \textbf{Fake}: implementacion alternativa funcional pero simplificada (e.g.\ una base
  de datos en memoria en vez de la real). Suficiente para correr tests pero no para
  produccion.
  \item \textbf{Mock}: doble pre-programado con expectativas exactas sobre que metodos se le
  van a llamar, con que argumentos y en que orden. Si las expectativas no se cumplen, el
  test falla. Frameworks como \texttt{EasyMock} y \texttt{Mockito} lo automatizan.
\end{enumerate}

\begin{ejemplobox}
Con \texttt{EasyMock}:
\begin{lstlisting}[style=codestyle,language=Java]
LoginService mock = createMock(LoginService.class);
expect(mock.login("ip", "user", "hash")).andReturn(false).times(3);
replay(mock);
// ... ejercitar el SUT que usa el mock ...
verify(mock); // falla si las invocaciones no se cumplieron
\end{lstlisting}
\end{ejemplobox}

\subsection{Cuando usar cada uno}

\begin{itemize}
  \item \textbf{Dummy}: cuando el parametro no influye en el escenario.
  \item \textbf{Stub}: cuando solo se necesita una respuesta canned.
  \item \textbf{Spy}: cuando ademas se quiere verificar que el SUT llamo correctamente.
  \item \textbf{Fake}: cuando la integracion real es demasiado pesada pero hay logica
  significativa a ejercitar.
  \item \textbf{Mock}: cuando la \emph{interaccion} entre SUT y DOC es la propiedad central
  que se quiere verificar (e.g.\ ``el SUT debe llamar a \texttt{save} exactamente una vez'').
\end{itemize}

\begin{ideabox}
Regla practica: usar el doble \emph{mas simple} que sirva. Empezar con \texttt{stub}; subir a
\texttt{mock} solo cuando hace falta verificar interacciones especificas.
\end{ideabox}

\section{Generacion automatica de tests}

\subsection{Idea general}

Crear tests a mano es \emph{costoso}, \emph{tedioso} e \emph{incompleto}. La generacion
automatica busca \emph{producir tests} a partir del codigo fuente o de su firma, sin
intervencion humana en cada caso. Hay varias estrategias; el TP8 cubre dos representativas:

\begin{itemize}
  \item \textbf{Randoop}: generacion aleatoria guiada por feedback.
  \item \textbf{EvoSuite}: busqueda evolutiva guiada por cobertura.
\end{itemize}

\subsection{Randoop: generacion aleatoria guiada por feedback}

\paragraph{Algoritmo basico (simplificado).}

\begin{enumerate}
  \item Mantener un \emph{pool} de secuencias previas que terminan en objetos validos.
  \item Mientras quede tiempo:
  \begin{enumerate}
    \item Elegir un metodo $m(T_1, \dots, T_k) \to T_{\text{ret}}$ al azar de la API.
    \item Para cada parametro de tipo $T_i$, elegir aleatoriamente del pool una secuencia
    $S_i$ que construya un valor de ese tipo.
    \item Construir la nueva secuencia $S_{\text{new}} = S_1; S_2; \dots; S_k;\
    T_{\text{ret}}\ v_{\text{new}} = m(v_1, \dots, v_k);$.
    \item Si la secuencia ya existia sintacticamente, descartar.
    \item Ejecutar y \emph{clasificar} la secuencia:
    \begin{itemize}
      \item Crash inesperado o violacion de contrato basico $\to$ \texttt{ErrorTest}
      (bug a investigar).
      \item Ejecucion exitosa $\to$ \texttt{RegressionTest} (suite de regresion) y se agrega
      al pool.
      \item Excepcion declarada o invalidacion $\to$ descartar.
    \end{itemize}
  \end{enumerate}
\end{enumerate}

\paragraph{Salida.}
\begin{itemize}
  \item \texttt{ErrorTest}: tests que rompen contratos basicos (e.g.\
  \texttt{equals/hashCode}, \texttt{toString} con NPE).
  \item \texttt{RegressionTest}: tests que documentan el comportamiento actual.
\end{itemize}

\paragraph{Aserciones de regresion.} Randoop usa metodos \emph{observadores} (sin
side-effects) para capturar el estado actual del objeto y generar asserts del estilo
\texttt{assertTrue(b2 == false)}. Estos asserts no garantizan correctitud, pero detectan
\emph{regresiones} ante cambios futuros.

\paragraph{Mecanismos para mejorar.}

\begin{itemize}
  \item \texttt{--junit-before-each}/\texttt{--junit-after-each}: inyectar setup/teardown
  (e.g.\ recrear \texttt{System.in}, archivos temporales).
  \item \texttt{--system-props}: fijar locale, timezone para reproducibilidad.
  \item \texttt{--methodlist}/\texttt{--omitmethods}: concentrar la exploracion.
  \item \texttt{--timeout}: cortar bloqueos en lecturas de \texttt{stdin}.
  \item \texttt{@randoop.CheckRep}: indicar que un metodo es un invariante;
  Randoop lo evalua en cada secuencia.
\end{itemize}

\subsection{EvoSuite: busqueda evolutiva}

\paragraph{Idea.} En lugar de generar al azar, EvoSuite mantiene una \emph{poblacion} de
suites de test y la evoluciona usando un \emph{algoritmo genetico}:

\begin{itemize}
  \item \textbf{Funcion de fitness}: usualmente cobertura (lineas, ramas, \emph{weak
  mutation}). Cuanto mas alta, mejor.
  \item \textbf{Cruce y mutacion}: combina suites existentes y aplica pequenas mutaciones
  (agregar una llamada, cambiar un valor, eliminar una linea).
  \item \textbf{Seleccion}: las suites con mejor fitness sobreviven a la siguiente
  generacion.
\end{itemize}

\paragraph{Salida.}
\begin{itemize}
  \item Suite ``feliz'' (e.g.\ \texttt{Server\_ESTest.java}) con tests que pasan.
  \item Suite ``failed'' (e.g.\ \texttt{Server\_Failed\_ESTest.java}) con escenarios que
  rompen el estado interno.
\end{itemize}

\paragraph{Diferencia clave con Randoop.}

\begin{center}
\begin{tabular}{lll}
\toprule
& Randoop & EvoSuite \\
\midrule
Tecnica & aleatoria + feedback & busqueda evolutiva \\
Objetivo & inputs validos / errores & maximizar cobertura \\
Tiempo de generacion & muy rapido & mas lento (search budget) \\
Calidad de la suite & superficial & mas profunda \\
Costo de integracion & bajo (Java moderno) & alto (Java 8/11) \\
\bottomrule
\end{tabular}
\end{center}

\begin{ideabox}
\textbf{Cuando elegir cada uno.} Randoop es ideal como primera pasada: rapido, encuentra
crashes evidentes y NPEs. EvoSuite vale la pena cuando se busca cobertura alta de ramas; es
mas caro de configurar pero produce suites mas exigentes.
\end{ideabox}

\subsection{Limitaciones comunes}

\begin{itemize}
  \item \textbf{Oraculo limitado}. Sin propiedades explicitas, las herramientas solo detectan
  bugs ``visibles'': crashes, violaciones de contratos del lenguaje.
  \item \textbf{Tests poco legibles}. Las secuencias generadas son mecanicas y dificiles de
  interpretar.
  \item \textbf{Dependencias externas}. Cualquier clase que dependa de \texttt{System.in},
  filesystem, red o reloj suele degradar drasticamente la calidad de la generacion.
\end{itemize}

\section{Ejecucion simbolica (Symbolic Execution)}

\subsection{Idea}

\begin{definicionbox}
\textbf{Ejecucion simbolica} (SE). Es una tecnica que ejecuta el programa con \emph{valores
simbolicos} en lugar de concretos. En cada decision se bifurca: una rama asume que la
condicion vale \texttt{true}, otra que vale \texttt{false}. En cada paso se acumulan dos
piezas de estado:
\begin{itemize}
  \item \textbf{Estado simbolico}: el valor (formula) de cada variable.
  \item \textbf{Condicion de camino} (\emph{path condition}): la conjuncion de todas las
  guardas que se asumieron hasta llegar a ese punto.
\end{itemize}
\end{definicionbox}

Cada hoja del arbol simbolico corresponde a un camino del programa. Para fabricar un input
concreto que recorra ese camino, se le pasa la \emph{path condition} a un \textbf{constraint
solver} (e.g.\ Z3, CVC5).

\subsection{Ejemplo}

\begin{lstlisting}[style=codestyle,language=Java]
public int testMe(int a, int b) {
    if (b < 0) {
        return -a;
    } else {
        int c = a - b;
        if (b + c == 400) {
            throw new IllegalArgumentException();
        } else {
            return c;
        }
    }
}
\end{lstlisting}

Los tres caminos producen tres \emph{path conditions}:

\begin{enumerate}
  \item $b_0 < 0$ \quad (camino \texttt{return -a}).
  \item $b_0 \ge 0 \wedge b_0 + (a_0 - b_0) \ne 400$, simplificada $b_0 \ge 0 \wedge a_0 \ne
  400$ \quad (camino \texttt{return c}).
  \item $b_0 \ge 0 \wedge a_0 = 400$ \quad (camino \texttt{throw}).
\end{enumerate}

Pasados al constraint solver, dan ejemplos concretos:

\begin{itemize}
  \item $b_0 = -1,\ a_0 = 0$ (cualquiera).
  \item $b_0 = 0,\ a_0 = 0$.
  \item $b_0 = 0,\ a_0 = 400$.
\end{itemize}

\begin{ideabox}
La ejecucion simbolica es \emph{precisa}: cada test generado \emph{garantiza} recorrer un
camino especifico. Es lo opuesto a la generacion aleatoria, que puede ejecutar muchas veces
el mismo camino y nunca tocar otros.
\end{ideabox}

\subsection{Ejecucion simbolica \emph{dinamica} (concolic execution)}

La SE pura tiene problemas en presencia de codigo que no se puede analizar simbolicamente
(llamadas a librerias nativas, IO, codigo binario). La variante \emph{concolic} (concrete +
symbolic) combina ambas:

\begin{enumerate}
  \item Se elige un input concreto inicial.
  \item Se ejecuta el programa concretamente, pero \emph{en paralelo} se mantiene la
  ejecucion simbolica.
  \item Cuando aparece codigo no analizable, se usa el valor concreto.
  \item Al terminar, se invierte una guarda del camino y se le pide al solver un input que
  recorra el camino alternativo.
  \item Se repite, ganando un camino nuevo en cada iteracion.
\end{enumerate}

Herramientas tipicas: \emph{KLEE}, \emph{Symbolic PathFinder}, \emph{angr}, \emph{DART},
\emph{CUTE}.

\subsection{Constraint solvers}

Un \textbf{constraint solver} recibe una formula logica (tipicamente en logica de primer
orden sobre enteros, bitvectors, arrays, etc.) y responde:

\begin{itemize}
  \item \textbf{SAT}: la formula es satisfacible y devuelve un modelo concreto.
  \item \textbf{UNSAT}: no hay valores que la satisfagan; el camino es \emph{infactible}.
  \item \textbf{UNKNOWN}: el solver no pudo decidir en el tiempo dado.
\end{itemize}

Los \emph{SMT solvers} modernos (\emph{Satisfiability Modulo Theories}) como Z3 manejan
combinaciones de teorias (aritmetica, bitvectors, strings, arrays).

\subsection{Limitaciones de SE}

\begin{itemize}
  \item \textbf{Path explosion}: el numero de caminos crece exponencialmente con el numero de
  ramas. Bucles y recursion son particularmente problematicos.
  \item \textbf{Constraints no decidibles}: el solver puede fallar con aritmetica no lineal,
  strings complejos o llamadas externas.
  \item \textbf{Modelado de heap}: estructuras dinamicas (listas enlazadas, arboles) son
  dificiles de simbolizar de manera completa.
  \item \textbf{Codigo nativo}: las llamadas a JNI o librerias C escapan al analisis.
\end{itemize}

\section{Relacion entre las tres tecnicas}

\begin{center}
\begin{tabular}{lll}
\toprule
Tecnica & Que automatiza & Costo \\
\midrule
Mocks & aislar el SUT & escribir el doble \\
Randoop & generar tests aleatorios & casi nulo \\
EvoSuite & maximizar cobertura & search budget \\
SE / concolic & explorar caminos especificos & solver + analisis \\
\bottomrule
\end{tabular}
\end{center}

\begin{ideabox}
Las tecnicas no compiten: se combinan.
\begin{itemize}
  \item Mocks se usan \emph{dentro} de cualquier tipo de test (unit, PBT, generado).
  \item Randoop y EvoSuite suelen \emph{integrar} dobles automaticamente cuando detectan
  dependencias.
  \item La ejecucion simbolica puede usarse para \emph{guiar} la generacion (e.g.\ ayudando
  a EvoSuite a encontrar caminos profundos).
\end{itemize}
\end{ideabox}

\section{Comparacion con criterios anteriores}

\begin{itemize}
  \item \textbf{Particion del input} (TP2) y \textbf{PBT} (TP7) producen tests
  \emph{escritos a mano} pero \emph{abstractos}; las tecnicas del TP8 los generan
  automaticamente, con menos garantia semantica.
  \item \textbf{Cobertura estructural} (TP3-TP4) provee la metrica que la busqueda
  evolutiva de EvoSuite optimiza.
  \item \textbf{Cobertura logica} (TP5) y \textbf{mutacion} (TP6) son maneras de
  \emph{evaluar} una suite generada automaticamente.
  \item \textbf{Fuzzing} (TP6) es el primo pobre de la ejecucion simbolica: tambien busca
  inputs, pero al azar.
\end{itemize}

\section{Glosario}

\begin{itemize}
  \item \textbf{SUT} (\emph{System Under Test}): la pieza de software que se quiere
  testear.
  \item \textbf{DOC} (\emph{Depended-On Component}): componente del que depende el SUT.
  \item \textbf{Dummy/Stub/Spy/Fake/Mock}: tipos de dobles ordenados de menor a mayor
  funcionalidad.
  \item \textbf{Path condition}: conjuncion de las guardas asumidas en un camino simbolico.
  \item \textbf{Constraint solver}: herramienta que decide la satisfacibilidad de formulas
  logicas y produce modelos.
  \item \textbf{SAT/UNSAT/UNKNOWN}: posibles respuestas de un constraint solver.
  \item \textbf{Path explosion}: explosion combinatoria del numero de caminos en SE.
  \item \textbf{Search budget}: presupuesto de tiempo o de evaluaciones de fitness en una
  busqueda evolutiva.
  \item \textbf{ErrorTest / RegressionTest}: dos clases de tests que produce Randoop.
\end{itemize}

\section*{Conclusion}
\addcontentsline{toc}{section}{Conclusion}

Los dobles de prueba resuelven el problema de \emph{como} testear unidades acopladas. La
generacion automatica de tests (Randoop, EvoSuite) resuelve el problema de \emph{escribir}
tests cuando hay mucho codigo. La ejecucion simbolica resuelve el problema de
\emph{recorrer} caminos especificos del programa con precision. Ninguna de las tres
reemplaza al criterio humano: el oraculo --- saber que es ``correcto'' --- sigue siendo
responsabilidad del desarrollador. Pero las tres juntas permiten escalar el testing mucho
mas alla de lo que un humano podria escribir caso por caso, sin perder el control sobre
\emph{que} se esta probando.

\end{document}
